
\subsection{Objetivo}
Una cuestión planteada en las reuniones se relaciona con el modelo de negocio de YPF Boxes, el cual actualmente opera con capacidad ociosa. El objetivo de este módulo es desarrollar una estrategia que permita brindar información útil sobre las oportunidades y desafíos que presenta este modelo de negocios. La primera pregunta que se propone abordar se refiere a la definición y comprensión del marco en el que opera YPF Boxes, en particular su entorno competitivo y la definición del mercado geográfico relevante. 

Sin embargo, la identificación del mercado geográfico relevante en este caso no resulta directa. A diferencia de otros servicios ofrecidos en estaciones de servicio, YPF Boxes compite no solo con servicios equivalentes dentro de otras estaciones, sino también con establecimientos especializados como lubricentros independientes. Estos competidores pueden presentar características heterogéneas, tales como distintos niveles de antigüedad en el mercado, clientelas previamente establecidas y grados variables de reconocimiento local. Asimismo, la demanda por servicios asociados a YPF Boxes puede exhibir patrones de estacionalidad y picos de utilización en determinados momentos del tiempo, lo que introduce variación adicional en la intensidad competitiva observada. En conjunto, estos factores dificultan una delimitación simple del mercado relevante y refuerzan la necesidad de realizar un análisis empírico sistemático que permita caracterizar adecuadamente el entorno competitivo en el que operan estos servicios.

\subsection{Diseño econométrico}
El presente bloque propone estimar un modelo econométrico cuyo objetivo es analizar los determinantes del grado de utilización de los YPF Boxes en estaciones de servicio. En particular, se busca cuantificar cómo distintos factores relacionados con la competencia local, las características del punto de venta y variables temporales inciden sobre la intensidad de uso de estas instalaciones. El análisis permitirá identificar patrones sistemáticos de demanda y contribuir a comprender en qué contextos operativos los YPF Boxes alcanzan mayores niveles de utilización.

La variable dependiente del modelo será el grado de utilización del YPF Box en cada estación y período de tiempo considerado. Esta variable puede definirse como la proporción del tiempo disponible en el que los boxes se encuentran ocupados, o bien como una medida equivalente de intensidad de uso construida a partir de registros operativos. La especificación empírica se estimará a nivel estación–período, lo que permitirá explotar variación tanto transversal como temporal en los determinantes del uso de la infraestructura.

Entre las variables explicativas principales se incorporará una medida de competencia local, capturada a través de la cantidad de estaciones competidoras ubicadas en proximidad geográfica a cada estación de YPF. Asimismo, el modelo incluirá una variable indicadora que distinga si la estación se encuentra ubicada sobre una ruta o dentro de una zona urbana, con el objetivo de capturar diferencias estructurales en los patrones de circulación y en la demanda potencial de servicios.

Adicionalmente, se incluirán variables relacionadas con las condiciones comerciales y operativas de cada estación. En particular, se considerará el precio de los productos o servicios ofrecidos, así como la antigüedad del local, entendida como el tiempo transcurrido desde la apertura o desde la implementación de la infraestructura relevante. Estas variables permiten capturar posibles efectos asociados a estrategias de precios, posicionamiento comercial y maduración del establecimiento.

Finalmente, el modelo incorporará controles temporales que permitan capturar patrones sistemáticos de estacionalidad y variaciones de corto plazo en la demanda. En particular, se incluirán efectos correspondientes al mes del año y a la semana, con el fin de controlar por fluctuaciones asociadas a cambios estacionales, ciclos de actividad y patrones recurrentes de utilización. Esta especificación permitirá aislar con mayor precisión el efecto de las variables estructurales de interés sobre el grado de utilización de los YPF Boxes.

\subsection{Datos requeridos}

Para la implementación del análisis propuesto será necesario contar, en primer lugar, con información operativa detallada sobre los YPF Boxes. En particular, se requerirán datos que permitan medir la capacidad instalada y el grado de utilización de los boxes a nivel de estación y período de tiempo. Asimismo, resultará relevante disponer de información sobre las características de cada estación, incluyendo su localización geográfica, si se encuentra ubicada en ruta o en zona urbana, la antigüedad del establecimiento, y los precios de los productos o servicios asociados. Esta información permitirá construir la variable dependiente del modelo y caracterizar las condiciones comerciales y operativas bajo las cuales se ofrece el servicio.

Adicionalmente, será necesario incorporar datos provenientes de fuentes externas que permitan caracterizar el entorno competitivo de cada estación. En particular, se requerirá información georreferenciada sobre la ubicación de establecimientos competidores, incluyendo tanto estaciones de servicio que ofrezcan servicios similares como lubricentros independientes, con el fin de construir indicadores de densidad competitiva y proximidad geográfica. De manera complementaria, se utilizarán variables temporales —como mes del año y semana— que permitan capturar patrones de estacionalidad y variaciones en la demanda. La integración de estas distintas fuentes de información permitirá construir una base de datos adecuada para la estimación del modelo econométrico propuesto.

